
\documentclass[11pt,a4paper]{article}
\usepackage[margin=0.9in,headheight=14pt]{geometry}
\usepackage[T1]{fontenc}
\usepackage{lmodern}
\usepackage{amsmath,amssymb}
\usepackage{enumitem,booktabs,fancyhdr}
\setlength{\parindent}{0pt}
\setlength{\parskip}{0.5em}
\linespread{1.08}
\pagestyle{fancy}
\fancyhf{}
\fancyhead[C]{\small Probability \& Statistics for EECS: Homework \#04}
\fancyfoot[C]{\thepage}
\renewcommand{\headrulewidth}{0.4pt}
\renewcommand{\footrulewidth}{0.4pt}
\setlist[enumerate]{label=(\alph*),leftmargin=*,itemsep=1em}
\newcommand{\Prob}{P}
\title{SI 140A.02 Probability \& Statistics for EECS, Fall 2026\\[0.4em]Homework 4}
\author{Name: \underline{\hspace{2.2in}}\\[0.8em]Student ID: \underline{\hspace{1.9in}}}
\date{
Due date: Due on Oct. 18, 2026, 23:59 UTC+8}
\begin{document}
\maketitle
\thispagestyle{plain}
\vspace{1em}
There are five problems in this assignment.

Read all the instructions below carefully before you start working on the assignment, and before you make a submission.
\begin{itemize}
\item You are required to write down all the major steps towards making your conclusions; otherwise you may receive limited credit for the problem.
\item Write your homework in English; otherwise you will receive no credit for this homework.
\item Any form of plagiarism will result in no credit for this homework.
\item You are required to submit your homework in PDF format generated by LaTeX, otherwise you will receive no credit for this homework.
\end{itemize}


\clearpage
\section*{Problem 1}

A token starts at position $1$ on the four positions $0,1,2,3$.
It stops when it first reaches $0$ or $3$.
Fix $r\in[0,1]$.
All random choices in different rounds are independent.
In an interior position, either of the following rules can be used
throughout the process:

\begin{itemize}
    \item \textbf{Pause rule:}
    With probability $r$, the token stays where it is;
    with probability $1-r$, a fresh fair coin moves it
    one position left or right.

    \item \textbf{Restart rule:}
    With probability $r$, the token is placed at position $1$;
    with probability $1-r$, a fresh fair coin moves it
    one position left or right.
\end{itemize}

\begin{enumerate}
    \item Under the pause rule, determine the probability
    of ever reaching position $3$ for every $r\in[0,1]$.
    Justify your answer.

    \item Answer part (a) under the restart rule.

    \item A student claims that pausing and restarting
    have the same effect because both only delay the token.
    Determine exactly when the two success probabilities are equal.
    For $r<1$, let $T$ be the number of rounds until the token stops.
    Find $\Pr(T>n)$ for every integer $n\geq0$ under each rule,
    and use your result to prove that each process stops
    with probability $1$.
    \emph{Hint:} Condition on either possible interior position
    after $n$ rounds.
\end{enumerate}

\subsection*{Solution}

\begin{enumerate}

%---------------- Part (a) ----------------
\item[(a)]

Let $p_i$ denote the probability of reaching position $3$
before position $0$, starting from position $i$.

\[
p_0=0,\qquad p_3=1.
\]

By conditioning on the first round,

\[
\begin{cases}
p_1=rp_1+\dfrac{1-r}{2}p_0
          +\dfrac{1-r}{2}p_2,\\[8pt]
p_2=rp_2+\dfrac{1-r}{2}p_1
          +\dfrac{1-r}{2}p_3.
\end{cases}
\]

For $0\le r<1$,

\[
\begin{cases}
(1-r)p_1=\dfrac{1-r}{2}p_2,\\[8pt]
(1-r)p_2=\dfrac{1-r}{2}(p_1+1).
\end{cases}
\]

\[
\Rightarrow
\begin{cases}
p_1=\dfrac12p_2,\\[6pt]
p_2=\dfrac12(p_1+1).
\end{cases}
\]

\[
\Rightarrow
2p_1=\frac12(p_1+1)
\Rightarrow 4p_1=p_1+1
\Rightarrow p_1=\frac13.
\]

For $r=1$, the token remains at position $1$ forever,

\[
p_1=0.
\]

Thus,

\[
\Pr(\text{ever reaching }3)=
\begin{cases}
\dfrac13, & 0\le r<1,\\[6pt]
0, & r=1.
\end{cases}
\]

%---------------- Part (b) ----------------
\item[(b)]

Let $p_i$ denote the probability of reaching position $3$
before position $0$ under the restart rule.

\[
p_0=0,\qquad p_3=1.
\]

By conditioning on the first round,

\[
\begin{cases}
p_1=rp_1+\dfrac{1-r}{2}p_0
          +\dfrac{1-r}{2}p_2,\\[8pt]
p_2=rp_1+\dfrac{1-r}{2}p_1
          +\dfrac{1-r}{2}p_3.
\end{cases}
\]

For $0\le r<1$,

\[
\begin{cases}
p_1=\dfrac12p_2,\\[8pt]
p_2=\dfrac{1+r}{2}p_1+\dfrac{1-r}{2}.
\end{cases}
\]

\[
\Rightarrow
2p_1=\frac{1+r}{2}p_1+\frac{1-r}{2}
\]

\[
\Rightarrow
4p_1=(1+r)p_1+1-r
\]

\[
\Rightarrow
(3-r)p_1=1-r
\]

\[
\Rightarrow
p_1=\frac{1-r}{3-r}.
\]

For $r=1$, the token remains at position $1$ forever,

\[
p_1=0=\frac{1-r}{3-r}.
\]

Thus,

\[
\Pr(\text{ever reaching }3)
=\frac{1-r}{3-r},
\qquad 0\le r\le1.
\]

%---------------- Part (c) ----------------
\item[(c)]

From parts (a) and (b),

\[
P_{\mathrm{pause}}=
\begin{cases}
\dfrac13, & 0\le r<1,\\[6pt]
0, & r=1,
\end{cases}
\qquad
P_{\mathrm{restart}}=\frac{1-r}{3-r}.
\]

For $0\le r<1$,

\[
P_{\mathrm{pause}}=P_{\mathrm{restart}}
\]

\[
\Leftrightarrow
\frac13=\frac{1-r}{3-r}
\]

\[
\Leftrightarrow
3-r=3-3r
\Leftrightarrow r=0.
\]

For $r=1$,

\[
P_{\mathrm{pause}}=P_{\mathrm{restart}}=0.
\]

Thus, the two success probabilities are equal exactly when

\[
r\in\{0,1\}.
\]

\medskip
\noindent
\textbf{Distribution of the stopping time.}

For $0\le r<1$, let $X_n$ denote the position after
$n$ rounds, and define

\[
S_n=\Pr(T>n),\qquad S_0=1.
\]

Under the pause rule, for $i=1,2$,

\[
\Pr(T>n+1\mid T>n,X_n=i)
=r+\frac{1-r}{2}
=\frac{1+r}{2}.
\]

By the law of total probability,

\[
\begin{aligned}
S_{n+1}
&=\sum_{i=1}^{2}
\Pr(T>n,X_n=i)\\
&\qquad\cdot
\Pr(T>n+1\mid T>n,X_n=i)\\[4pt]
&=\frac{1+r}{2}
\sum_{i=1}^{2}\Pr(T>n,X_n=i)\\[4pt]
&=\frac{1+r}{2}S_n.
\end{aligned}
\]

\[
\because S_0=1
\]

\[
\Rightarrow
S_n=\left(\frac{1+r}{2}\right)^n.
\]

Under the restart rule, for $i=1,2$,

\[
\Pr(T>n+1\mid T>n,X_n=i)
=r+\frac{1-r}{2}
=\frac{1+r}{2}.
\]

Therefore, the same recurrence holds:

\[
S_{n+1}=\frac{1+r}{2}S_n,
\qquad S_0=1.
\]

\[
\Rightarrow
S_n=\left(\frac{1+r}{2}\right)^n.
\]

Thus, under both rules,

\[
\Pr(T>n)
=\left(\frac{1+r}{2}\right)^n,
\qquad n\ge0,\quad 0\le r<1.
\]

\medskip
\noindent
\textbf{Almost sure termination.}

\[
\because\quad
0\le r<1
\Rightarrow
0<\frac{1+r}{2}<1
\]

\[
\Rightarrow
\lim_{n\to\infty}\Pr(T>n)
=
\lim_{n\to\infty}
\left(\frac{1+r}{2}\right)^n
=0.
\]

\[
\Pr(T=\infty)
=\lim_{n\to\infty}\Pr(T>n)=0.
\]

\[
\Rightarrow
\Pr(T<\infty)
=1-\Pr(T=\infty)
=1.
\]

Thus, both processes stop with probability $1$
for every $0\le r<1$.

\end{enumerate}


\clearpage
\section*{Problem 2}

Two algorithms, $A$ and $B$, are each tested on $100$ tasks.
Across the $200$ test records, exactly $100$ tasks are easy
and $100$ are hard. Their exact success rates are:

\begin{center}
\begin{tabular}{ccc}
\toprule
Algorithm & Easy tasks & Hard tasks\\
\midrule
$A$ & $4/5$ & $2/5$\\
$B$ & $9/10$ & $3/5$\\
\bottomrule
\end{tabular}
\end{center}

Each algorithm's easy/hard versus success/failure table
has strictly positive integer entries. Let $x$ be the
number of easy tasks assigned to $A$.

\begin{enumerate}
\item Find every feasible $x$ for which $A$ has a strictly
greater overall success rate than $B$. For the smallest
such $x$, give both complete $2\times2$ count tables.

\textbf{Solution:}

The numbers of easy and hard tasks are
\[
\begin{array}{c|cc}
 & \text{Easy} & \text{Hard}\\
\hline
A & x & 100-x\\
B & 100-x & x
\end{array}
\]

All success and failure counts must be strictly
positive integers:
\[
\begin{aligned}
&5\mid x,\qquad 5\mid(100-x),\\
&10\mid(100-x),\qquad 5\mid x.
\end{aligned}
\]

\[
\implies 10\mid x,\qquad 0<x<100
\]

\[
\therefore\quad
x\in\{10,20,30,40,50,60,70,80,90\}.
\]

The total numbers of successes are
\[
\begin{aligned}
S_A
&=\frac45x+\frac25(100-x)\\
&=40+\frac25x,\\[6pt]
S_B
&=\frac9{10}(100-x)+\frac35x\\
&=90-\frac3{10}x.
\end{aligned}
\]

\[
\begin{aligned}
P(S\mid A)>P(S\mid B)
&\iff S_A>S_B\\
&\iff 40+\frac25x>90-\frac3{10}x\\
&\iff \frac7{10}x>50\\
&\iff x>\frac{500}{7}.
\end{aligned}
\]

Combining the feasible values,
\[
x\in\{80,90\}.
\]

For the smallest value $x=80$:

\[
\text{Algorithm }A:
\]

\begin{center}
\begin{tabular}{c|cc}
\toprule
 & Success & Failure\\
\midrule
Easy & 64 & 16\\
Hard & 8 & 12\\
\bottomrule
\end{tabular}
\end{center}

\[
\text{Algorithm }B:
\]

\begin{center}
\begin{tabular}{c|cc}
\toprule
 & Success & Failure\\
\midrule
Easy & 18 & 2\\
Hard & 48 & 32\\
\bottomrule
\end{tabular}
\end{center}

\[
\begin{aligned}
P(S\mid A)&=\frac{64+8}{100}=\frac{18}{25},\\
P(S\mid B)&=\frac{18+48}{100}=\frac{33}{50}.
\end{aligned}
\]

\[
\therefore\quad
P(S\mid A)=0.72>0.66=P(S\mid B).
\]


\item From this part onward, use a probability model
without finite-sample or integer-count constraints.
The displayed within-type success probabilities remain
fixed. Both algorithms receive the same easy-task
proportion $t\in[0,1]$. An audit reports that their
overall success probabilities differ by $13/100$,
but omits which algorithm performs better.
Recover $t$, identify the better algorithm, and
calculate both overall success probabilities.

\textbf{Solution:}

By the law of total probability,
\[
\begin{aligned}
P(S\mid A)
&=\frac45t+\frac25(1-t)\\
&=\frac25+\frac25t,\\[6pt]
P(S\mid B)
&=\frac9{10}t+\frac35(1-t)\\
&=\frac35+\frac3{10}t.
\end{aligned}
\]

\[
\begin{aligned}
&|P(S\mid A)-P(S\mid B)|=\frac{13}{100}\\
\iff\;&
\left|\frac1{10}t-\frac15\right|
=\frac{13}{100}\\
\iff\;&
|10t-20|=13.
\end{aligned}
\]

\[
\implies
\begin{cases}
10t-20=13 \implies t=\frac{33}{10},\\
10t-20=-13 \implies t=\frac7{10}.
\end{cases}
\]

\[
\because\quad t\in[0,1]
\]

\[
\therefore\quad t=\frac7{10}.
\]

Substituting,
\[
\begin{aligned}
P(S\mid A)
&=\frac25+\frac25\cdot\frac7{10}\\
&=\frac{17}{25}=0.68,\\[6pt]
P(S\mid B)
&=\frac35+\frac3{10}\cdot\frac7{10}\\
&=\frac{81}{100}=0.81.
\end{aligned}
\]

\[
P(S\mid B)-P(S\mid A)
=\frac{81}{100}-\frac{68}{100}
=\frac{13}{100}.
\]

\[
\therefore\quad
P(S\mid B)>P(S\mid A).
\]


\item Now let $s,t\in(0,1)$ be the easy-task
proportions for $A,B$, respectively, so each
algorithm receives both task types. Construct one
pair $(s,t)$ satisfying $|s-t|\leq2/5$ for which
$A$'s overall success probability is exactly
$1/100$ greater than $B$'s. Calculate both
probabilities and check all the constraints for
your pair. Briefly explain how this is compatible
with $B$'s superiority within each task type.

\textbf{Solution:}

By the law of total probability,
\[
\begin{aligned}
P(S\mid A)
&=\frac45s+\frac25(1-s)\\
&=\frac25+\frac25s,\\[6pt]
P(S\mid B)
&=\frac9{10}t+\frac35(1-t)\\
&=\frac35+\frac3{10}t.
\end{aligned}
\]

The required conditions are
\[
\begin{cases}
P(S\mid A)-P(S\mid B)=\frac1{100},\\
|s-t|\leq\frac25,\\
0<s<1,\quad 0<t<1.
\end{cases}
\]

\[
\begin{aligned}
&\frac25+\frac25s
-\frac35-\frac3{10}t=\frac1{100}\\
\iff\;&40s-30t=21.
\end{aligned}
\]

Choose
\[
s-t=\frac25.
\]

\[
\begin{aligned}
40\left(t+\frac25\right)-30t&=21\\
10t+16&=21\\
t&=\frac12.
\end{aligned}
\]

\[
s=t+\frac25
=\frac12+\frac25
=\frac9{10}.
\]

\[
\therefore\quad
(s,t)=\left(\frac9{10},\frac12\right).
\]

The overall success probabilities are
\[
\begin{aligned}
P(S\mid A)
&=\frac45\cdot\frac9{10}
+\frac25\cdot\frac1{10}\\
&=\frac{19}{25}=0.76,\\[6pt]
P(S\mid B)
&=\frac9{10}\cdot\frac12
+\frac35\cdot\frac12\\
&=\frac34=0.75.
\end{aligned}
\]

Verification:
\[
\begin{aligned}
&0<\frac9{10}<1,\qquad
0<\frac12<1,\\[4pt]
&|s-t|
=\left|\frac9{10}-\frac12\right|
=\frac25\leq\frac25,\\[4pt]
&P(S\mid A)-P(S\mid B)\\
&\qquad=\frac{19}{25}-\frac34
=\frac1{100}.
\end{aligned}
\]

Although
\[
\frac9{10}>\frac45,
\qquad
\frac35>\frac25,
\]
algorithm $A$ receives a larger proportion of easy
tasks:
\[
s=\frac9{10}>\frac12=t.
\]

Thus, the different task-type weights reverse
the overall comparison, illustrating Simpson's
paradox.

\end{enumerate}


\clearpage
\section*{Problem 3}

An automated station is placed in mode $H$ with probability $1/3$
and mode $L$ with probability $2/3$. The mode is chosen once for
three checks. Let $A_i$ mean that check $i$ passes. Each check
passes with probability $3/4$ conditional on $H$, and with
probability $1/4$ conditional on $L$.
Let $U=A_1\cap A_2\cap A_3$.

\begin{enumerate}
\item
First assume the three pass events are mutually independent
conditional on either mode. Compute $\Prob(H\mid U)$.
Are $A_1$ and $A_2$ independent without conditioning on the mode?
Support your answer with probabilities, and explain the difference
between the two independence statements.

\item
For the rest of this problem, assume only that every pair of
pass events is independent conditional on each mode.
For a fixed mode, write $p$ for the common pass probability
and $z$ for the probability of three passes.
Complete the table below and justify why the outcomes in each
grouped row must have equal probabilities.
Write $1$ for a pass and $0$ for a failure.

\begin{center}
\begin{tabular}{cc}
\toprule
Outcome & Conditional probability of each outcome\\
\midrule
$111$ & $z$\\
$110,101,011$ & $\cdots$\\
$100,010,001$ & $\cdots$\\
$000$ & $\cdots$\\
\bottomrule
\end{tabular}
\end{center}

Find the full feasible intervals for
$h=\Prob(U\mid H)$ and $\ell=\Prob(U\mid L)$.
Explain why every pair $(h,\ell)$ in those intervals gives
a valid model with the specified marginals and pairwise
independence.

\item
Find the exact set of possible values of $\Prob(H\mid U)$
under the assumptions of part (b).
Prove both that no other value is possible and that every
value in your set is attainable.
Explain why the single answer in part (a) cannot be inferred
from pairwise independence alone.
\end{enumerate}

\noindent\textbf{Solution:}

\begin{enumerate}

%-------------------- (a) --------------------
\item[(a)]
By conditional mutual independence,
\[
\begin{aligned}
\Prob(U\mid H)
&=\left(\frac34\right)^3=\frac{27}{64},\\
\Prob(U\mid L)
&=\left(\frac14\right)^3=\frac1{64}.
\end{aligned}
\]

By Bayes' theorem,
\[
\begin{aligned}
\Prob(H\mid U)
&=\frac{\Prob(U\mid H)\Prob(H)}
{\Prob(U\mid H)\Prob(H)+\Prob(U\mid L)\Prob(L)}\\
&=\frac{\frac{27}{64}\cdot\frac13}
{\frac{27}{64}\cdot\frac13+\frac1{64}\cdot\frac23}\\
&=\frac{27}{29}.
\end{aligned}
\]

To examine unconditional independence,
\[
\begin{aligned}
\Prob(A_1)
&=\frac34\cdot\frac13+\frac14\cdot\frac23
=\frac5{12},\\
\Prob(A_2)&=\frac5{12}.
\end{aligned}
\]

Thus,
\[
\Prob(A_1)\Prob(A_2)=\frac{25}{144}.
\]

On the other hand,
\[
\begin{aligned}
\Prob(A_1\cap A_2)
&=\Prob(A_1\cap A_2\mid H)\Prob(H)\\
&\quad+\Prob(A_1\cap A_2\mid L)\Prob(L)\\
&=\left(\frac34\right)^2\frac13
+\left(\frac14\right)^2\frac23\\
&=\frac{11}{48}
=\frac{33}{144}.
\end{aligned}
\]

\[
\Prob(A_1\cap A_2)\ne\Prob(A_1)\Prob(A_2)
\]

\[
\implies A_1,\ A_2
\text{ are not unconditionally independent.}
\]

However, for $M\in\{H,L\}$,
\[
\Prob(A_1\cap A_2\mid M)
=\Prob(A_1\mid M)\Prob(A_2\mid M).
\]

Conditional independence holds for each fixed mode,
while unconditional independence fails because both
events depend on the same randomly selected mode.

%-------------------- (b) --------------------
\item[(b)]
Fix a mode $M\in\{H,L\}$.
Write all probabilities conditional on $M$, and let
\[
\Prob(A_i\mid M)=p,\qquad
\Prob(A_i\cap A_j\mid M)=p^2
\quad(i\ne j),
\]
\[
z=\Prob(A_1\cap A_2\cap A_3\mid M).
\]

For outcomes with exactly two passes,
\[
\begin{aligned}
\Prob(110\mid M)
&=\Prob(A_1\cap A_2\mid M)-z\\
&=p^2-z.
\end{aligned}
\]

Similarly,
\[
\Prob(110\mid M)
=\Prob(101\mid M)
=\Prob(011\mid M)
=p^2-z.
\]

For outcomes with exactly one pass,
\[
\begin{aligned}
\Prob(A_1\mid M)
={}&\Prob(100\mid M)+\Prob(101\mid M)\\
&+\Prob(110\mid M)+\Prob(111\mid M),
\end{aligned}
\]
\[
p=\Prob(100\mid M)+2(p^2-z)+z.
\]

\[
\implies \Prob(100\mid M)=p-2p^2+z.
\]

By the identical marginal probabilities and pairwise
intersection probabilities,
\[
\Prob(100\mid M)
=\Prob(010\mid M)
=\Prob(001\mid M)
=p-2p^2+z.
\]

Since the eight outcomes are disjoint and exhaustive,
\[
\begin{aligned}
\Prob(000\mid M)
&=1-z-3(p^2-z)-3(p-2p^2+z)\\
&=1-3p+3p^2-z.
\end{aligned}
\]

The completed table is
\[
\begin{array}{cc}
\toprule
\text{Outcome}
&\text{Conditional probability of each outcome}\\
\midrule
111 & z\\
110,101,011 & p^2-z\\
100,010,001 & p-2p^2+z\\
000 & 1-3p+3p^2-z\\
\bottomrule
\end{array}
\]

The equality within each grouped row follows
from the common marginals $p$ and the common
pairwise intersections $p^2$.

All outcome probabilities must be nonnegative:
\[
\begin{cases}
z\ge0,\\
p^2-z\ge0,\\
p-2p^2+z\ge0,\\
1-3p+3p^2-z\ge0.
\end{cases}
\]

\[
\implies
\max\{0,2p^2-p\}
\le z\le
\min\{p^2,1-3p+3p^2\}.
\]

For mode $H$, $p=\frac34$ and $z=h$:
\[
\begin{aligned}
h&\ge\max\left\{
0,2\left(\frac34\right)^2-\frac34
\right\}=\frac38,\\
h&\le\min\left\{
\left(\frac34\right)^2,
1-3\left(\frac34\right)
+3\left(\frac34\right)^2
\right\}=\frac7{16}.
\end{aligned}
\]

\[
\frac38\le h\le\frac7{16}.
\]

For mode $L$, $p=\frac14$ and $z=\ell$:
\[
\begin{aligned}
\ell&\ge\max\left\{
0,2\left(\frac14\right)^2-\frac14
\right\}=0,\\
\ell&\le\min\left\{
\left(\frac14\right)^2,
1-3\left(\frac14\right)
+3\left(\frac14\right)^2
\right\}=\frac1{16}.
\end{aligned}
\]

\[
0\le\ell\le\frac1{16}.
\]

To verify sufficiency, for any $z$ in its feasible
interval, all eight probabilities in the table
are nonnegative, and
\[
\begin{aligned}
&z+3(p^2-z)+3(p-2p^2+z)\\
&\qquad +(1-3p+3p^2-z)=1.
\end{aligned}
\]

Moreover,
\[
\begin{aligned}
\Prob(A_1\mid M)
&=z+2(p^2-z)+(p-2p^2+z)=p,\\
\Prob(A_1\cap A_2\mid M)
&=z+(p^2-z)=p^2.
\end{aligned}
\]

The same equalities hold for all $A_i$ and
all pairs $(A_i,A_j)$ by symmetry.

\[
\implies
\Prob(A_i\cap A_j\mid M)
=\Prob(A_i\mid M)\Prob(A_j\mid M).
\]

Thus every feasible $z$ defines a valid
conditional distribution with the required
marginals and pairwise independence.

The conditional distributions for $H$ and $L$
can be chosen independently, with
\[
\Prob(H)=\frac13,\qquad
\Prob(L)=\frac23.
\]

Consequently, every pair
\[
(h,\ell)\in
\left[\frac38,\frac7{16}\right]
\times\left[0,\frac1{16}\right]
\]
is attainable.

%-------------------- (c) --------------------
\item[(c)]
By Bayes' theorem,
\[
\begin{aligned}
\Prob(H\mid U)
&=\frac{\Prob(U\mid H)\Prob(H)}
{\Prob(U\mid H)\Prob(H)+\Prob(U\mid L)\Prob(L)}\\
&=\frac{\frac13h}{\frac13h+\frac23\ell}\\
&=\frac{h}{h+2\ell}
=\frac1{1+\frac{2\ell}{h}}.
\end{aligned}
\]

From part (b),
\[
\frac38\le h\le\frac7{16},
\qquad
0\le\ell\le\frac1{16}.
\]

\[
0\le\frac{2\ell}{h}
\le\frac{2\cdot\frac1{16}}{\frac38}
=\frac13.
\]

\[
\implies
\frac34\le\frac1{1+\frac{2\ell}{h}}\le1.
\]

Therefore,
\[
\frac34\le\Prob(H\mid U)\le1.
\]

The lower endpoint is attained at
\[
(h,\ell)=\left(\frac38,\frac1{16}\right),
\qquad
\Prob(H\mid U)=\frac34.
\]

The upper endpoint is attained whenever $\ell=0$:
\[
\Prob(H\mid U)=1.
\]

To prove that every intermediate value is attainable,
let
\[
t\in\left[\frac34,1\right].
\]

Choose
\[
h=\frac38,\qquad
\ell=\frac{3(1-t)}{16t}.
\]

Since $\frac34\le t\le1$,
\[
0\le\frac{3(1-t)}{16t}\le\frac1{16}.
\]

Thus $(h,\ell)$ is feasible by part (b).

Furthermore,
\[
\begin{aligned}
\Prob(H\mid U)
&=\frac{h}{h+2\ell}\\
&=\frac{\frac38}
{\frac38+2\cdot\frac{3(1-t)}{16t}}\\
&=\frac{\frac38}
{\frac38+\frac{3(1-t)}{8t}}\\
&=\frac{\frac38}{\frac{3}{8t}}\\
&=t.
\end{aligned}
\]

Hence every $t\in[\frac34,1]$ is attainable,
and no value outside this interval is possible.

The exact set is
\[
\left\{\Prob(H\mid U)\right\}
=\left[\frac34,1\right].
\]

Under conditional mutual independence in part (a),
\[
h=\left(\frac34\right)^3=\frac{27}{64},
\qquad
\ell=\left(\frac14\right)^3=\frac1{64},
\]
\[
\implies
\Prob(H\mid U)
=\frac{27}{27+2}
=\frac{27}{29}.
\]

However, conditional pairwise independence only gives
\[
\Prob(A_i\cap A_j\mid M)=p^2,
\]
without requiring
\[
\Prob(A_1\cap A_2\cap A_3\mid M)=p^3.
\]

Thus $h$ and $\ell$ are not uniquely determined,
and neither is $\Prob(H\mid U)$.

\end{enumerate}




\clearpage
\section*{Problem 4}

In a three-door game, the car is behind doors $1,2,3$ with probabilities
$1/2,1/3,1/6$. You choose door $1$. The host knows the car's location and
opens a goat door from doors $2,3$. If both doors contain goats, host A
opens door $2$ with probability $1/4$, whereas host B does so with
probability $3/4$. Otherwise, the opening is forced.

A fair coin selects one host, independently of the car. The selected host
first performs a calibration round, with the car known to be behind door
$1$ and door $1$ initially selected, and then the main game. Given the
host, the two rounds use independent random choices. Let $C$ and $J$
be the calibration and main openings, respectively, and let $K$ be the
main car door.

System F retains exactly the runs with $C=2$; system M retains exactly
the runs with $C=J$. Write $F=\{C=2\}$ and $M=\{C=J\}$. In a retained
run, the system and $J$ are reported, but the host and $C$ are not
separately reported. All questions below concern the main game.

\begin{enumerate}

\item Prove that the two systems give the same
posterior probability of $K=1$ when $J=2$ is reported, and calculate
this probability. Identify the conditioning events before calculating.

\item For each system, make a two-row table, indexed
by $j=2,3$, giving
\[
\Prob(J=j\mid\text{retained}),\qquad
\Prob(K=1\mid J=j,\text{retained}).
\]
Reuse your result from part (a). Apply the conditional law of total
probability to find \mbox{$\Prob(K=1\mid F)$} and
\mbox{$\Prob(K=1\mid M)$}.
Explain why the equality in part (a) does not imply equality of these
two probabilities.

\item Modify only the host selection: now select a
host independently by a fair coin for each of the two rounds.
Conditional on the selected hosts, the rounds use independent random
choices. Both recording systems keep their original rules. Does either
system change the joint distribution of $(K,J)$ by retaining runs?
Prove your answer for every pair $(k,j)$ with positive probability.
Begin by finding the probability of retention conditional on $K=k$
and $J=j$, using the law of total probability over the calibration
host. Explain what feature of the original experiment allowed the
two systems to give different answers in part (b).

\end{enumerate}

\noindent\textbf{Solution:}

Let $H\in\{A,B\}$ denote the host in the original experiment.

\[
\Prob(H=A)=\Prob(H=B)=\frac12,
\]
\[
\Prob(K=1)=\frac12,\quad
\Prob(K=2)=\frac13,\quad
\Prob(K=3)=\frac16.
\]

Define
\[
q_A=\frac14,\qquad q_B=\frac34.
\]

Then
\[
\Prob(C=2\mid H=h)=q_h,
\]
and
\[
\Prob(J=2\mid K=k,H=h)=
\begin{cases}
q_h, & k=1,\\
0, & k=2,\\
1, & k=3.
\end{cases}
\]

Given $H$, the two rounds are independent, and $H$ is independent
of $K$.

\medskip
\noindent\textbf{(a)}

The conditioning events are
\[
\begin{aligned}
F\cap\{J=2\}
&=\{C=2,J=2\},\\
M\cap\{J=2\}
&=\{C=J,J=2\}\\
&=\{C=2,J=2\}.
\end{aligned}
\]

Thus
\[
F\cap\{J=2\}=M\cap\{J=2\},
\]
\[
\implies
\Prob(K=1\mid F,J=2)
=
\Prob(K=1\mid M,J=2).
\]

By the law of total probability over $H$,
\[
\begin{aligned}
&\Prob(K=1,C=2,J=2)\\
&=\sum_{h\in\{A,B\}}
\Prob(H=h)\Prob(K=1)
\Prob(C=2\mid H=h)
\Prob(J=2\mid K=1,H=h)\\
&=\frac12\cdot\frac12
\left[
\left(\frac14\right)^2+
\left(\frac34\right)^2
\right]\\
&=\frac5{32}.
\end{aligned}
\]

For $K=2$,
\[
\Prob(K=2,C=2,J=2)=0.
\]

For $K=3$,
\[
\begin{aligned}
&\Prob(K=3,C=2,J=2)\\
&=\frac16\cdot\frac12
\left(\frac14+\frac34\right)\\
&=\frac1{12}.
\end{aligned}
\]

Therefore,
\[
\begin{aligned}
\Prob(C=2,J=2)
&=\sum_{k=1}^{3}\Prob(K=k,C=2,J=2)\\
&=\frac5{32}+0+\frac1{12}\\
&=\frac{23}{96}.
\end{aligned}
\]

Hence,
\[
\begin{aligned}
\Prob(K=1\mid C=2,J=2)
&=\frac{\Prob(K=1,C=2,J=2)}
{\Prob(C=2,J=2)}\\
&=\frac{5/32}{23/96}\\
&=\frac{15}{23}.
\end{aligned}
\]

\[
\implies
\Prob(K=1\mid F,J=2)
=
\Prob(K=1\mid M,J=2)
=
\frac{15}{23}.
\]

\medskip
\noindent\textbf{(b)}

\textbf{System F:}

\[
F=\{C=2\}.
\]

\[
\begin{aligned}
\Prob(F)
&=\frac12\cdot\frac14+
\frac12\cdot\frac34\\
&=\frac12.
\end{aligned}
\]

From part (a),
\[
\Prob(C=2,J=2)=\frac{23}{96}.
\]

Thus
\[
\begin{aligned}
\Prob(J=2\mid F)
&=\frac{\Prob(C=2,J=2)}{\Prob(F)}\\
&=\frac{23}{48},
\end{aligned}
\]
\[
\Prob(J=3\mid F)
=1-\frac{23}{48}
=\frac{25}{48}.
\]

For $K=1,C=2,J=3$,
\[
\begin{aligned}
&\Prob(K=1,C=2,J=3)\\
&=\frac12\cdot\frac12
\left[
\frac14\cdot\frac34+
\frac34\cdot\frac14
\right]\\
&=\frac3{32}.
\end{aligned}
\]

Also,
\[
\begin{aligned}
\Prob(C=2,J=3)
&=\Prob(C=2)-\Prob(C=2,J=2)\\
&=\frac12-\frac{23}{96}\\
&=\frac{25}{96}.
\end{aligned}
\]

Therefore,
\[
\begin{aligned}
\Prob(K=1\mid J=3,F)
&=\frac{\Prob(K=1,C=2,J=3)}
{\Prob(C=2,J=3)}\\
&=\frac{3/32}{25/96}\\
&=\frac9{25}.
\end{aligned}
\]

\[
\begin{array}{c|cc}
j & \Prob(J=j\mid F)
& \Prob(K=1\mid J=j,F)\\
\hline
2 & \dfrac{23}{48} & \dfrac{15}{23}\\[6pt]
3 & \dfrac{25}{48} & \dfrac9{25}
\end{array}
\]

\medskip
\textbf{System M:}

\[
M=\{C=J\}.
\]

\[
\begin{aligned}
&\Prob(K=1,C=3,J=3)\\
&=\frac12\cdot\frac12
\left[
\left(\frac34\right)^2+
\left(\frac14\right)^2
\right]\\
&=\frac5{32}.
\end{aligned}
\]

For $K=2$, the main opening is forced to $J=3$:
\[
\begin{aligned}
&\Prob(K=2,C=3,J=3)\\
&=\frac13\cdot\frac12
\left(\frac34+\frac14\right)\\
&=\frac16.
\end{aligned}
\]

For $K=3$,
\[
\Prob(K=3,C=3,J=3)=0.
\]

Thus
\[
\begin{aligned}
\Prob(C=3,J=3)
&=\frac5{32}+\frac16\\
&=\frac{31}{96}.
\end{aligned}
\]

\[
\begin{aligned}
\Prob(M)
&=\Prob(C=2,J=2)+\Prob(C=3,J=3)\\
&=\frac{23}{96}+\frac{31}{96}\\
&=\frac9{16}.
\end{aligned}
\]

\[
\begin{aligned}
\Prob(J=2\mid M)
&=\frac{23/96}{9/16}\\
&=\frac{23}{54},\\[6pt]
\Prob(J=3\mid M)
&=\frac{31/96}{9/16}\\
&=\frac{31}{54}.
\end{aligned}
\]

\[
\begin{aligned}
\Prob(K=1\mid J=3,M)
&=\frac{\Prob(K=1,C=3,J=3)}
{\Prob(C=3,J=3)}\\
&=\frac{5/32}{31/96}\\
&=\frac{15}{31}.
\end{aligned}
\]

\[
\begin{array}{c|cc}
j & \Prob(J=j\mid M)
& \Prob(K=1\mid J=j,M)\\
\hline
2 & \dfrac{23}{54} & \dfrac{15}{23}\\[6pt]
3 & \dfrac{31}{54} & \dfrac{15}{31}
\end{array}
\]

Applying the conditional law of total probability:

\[
\begin{aligned}
\Prob(K=1\mid F)
&=\sum_{j=2}^{3}
\Prob(K=1\mid J=j,F)\Prob(J=j\mid F)\\
&=\frac{15}{23}\cdot\frac{23}{48}
+\frac9{25}\cdot\frac{25}{48}\\
&=\frac{15}{48}+\frac9{48}\\
&=\frac12.
\end{aligned}
\]

Similarly,
\[
\begin{aligned}
\Prob(K=1\mid M)
&=\sum_{j=2}^{3}
\Prob(K=1\mid J=j,M)\Prob(J=j\mid M)\\
&=\frac{15}{23}\cdot\frac{23}{54}
+\frac{15}{31}\cdot\frac{31}{54}\\
&=\frac{15}{54}+\frac{15}{54}\\
&=\frac59.
\end{aligned}
\]

Hence,
\[
\Prob(K=1\mid F)=\frac12
\neq
\frac59=\Prob(K=1\mid M).
\]

Although
\[
\Prob(K=1\mid J=2,F)
=
\Prob(K=1\mid J=2,M),
\]
the probabilities for $J=3$ and the weights differ:
\[
\Prob(K=1\mid J=3,F)=\frac9{25}
\neq
\frac{15}{31}
=\Prob(K=1\mid J=3,M),
\]
\[
\Prob(J=2\mid F)=\frac{23}{48}
\neq
\frac{23}{54}=\Prob(J=2\mid M).
\]

Thus equality for $J=2$ does not imply equality after
averaging over $J$.

\medskip
\noindent\textbf{(c)}

Let $H_C$ and $H_J$ denote the independently selected hosts
for the calibration and main rounds.

\[
\Prob(H_C=A)=\Prob(H_C=B)=\frac12,
\]
\[
\Prob(H_J=A)=\Prob(H_J=B)=\frac12.
\]

The calibration host and its random choice are independent of
the main game.

For every $(k,j)$ with $\Prob(K=k,J=j)>0$, consider system F.

By the conditional law of total probability over $H_C$,
\[
\begin{aligned}
&\Prob(F\mid K=k,J=j)\\
&=\sum_{h\in\{A,B\}}
\Prob(H_C=h\mid K=k,J=j)\\
&\qquad\cdot
\Prob(C=2\mid H_C=h,K=k,J=j)\\
&=\frac12\cdot\frac14+
\frac12\cdot\frac34\\
&=\frac12.
\end{aligned}
\]

For system M,
\[
\begin{aligned}
&\Prob(M\mid K=k,J=j)\\
&=\Prob(C=j\mid K=k,J=j)\\
&=\sum_{h\in\{A,B\}}
\Prob(H_C=h\mid K=k,J=j)
\Prob(C=j\mid H_C=h).
\end{aligned}
\]

For $j=2$,
\[
\begin{aligned}
\Prob(M\mid K=k,J=2)
&=\frac12\cdot\frac14+
\frac12\cdot\frac34\\
&=\frac12.
\end{aligned}
\]

For $j=3$,
\[
\begin{aligned}
\Prob(M\mid K=k,J=3)
&=\frac12\cdot\frac34+
\frac12\cdot\frac14\\
&=\frac12.
\end{aligned}
\]

Thus, for every possible $(k,j)$,
\[
\Prob(F\mid K=k,J=j)
=
\Prob(M\mid K=k,J=j)
=
\frac12.
\]

The positive-probability pairs are
\[
(k,j)\in\{(1,2),(1,3),(2,3),(3,2)\}.
\]

Their joint probabilities are
\[
\begin{array}{c|cccc}
(k,j)&(1,2)&(1,3)&(2,3)&(3,2)\\
\hline
\Prob(K=k,J=j)
&\dfrac14&\dfrac14&\dfrac13&\dfrac16\\[6pt]
\Prob(F\mid K=k,J=j)
&\dfrac12&\dfrac12&\dfrac12&\dfrac12\\[6pt]
\Prob(M\mid K=k,J=j)
&\dfrac12&\dfrac12&\dfrac12&\dfrac12
\end{array}
\]

By the law of total probability,
\[
\begin{aligned}
\Prob(F)
&=\sum_{k,j}
\Prob(F\mid K=k,J=j)\Prob(K=k,J=j)\\
&=\frac12\sum_{k,j}\Prob(K=k,J=j)\\
&=\frac12.
\end{aligned}
\]

Similarly,
\[
\Prob(M)=\frac12.
\]

Applying Bayes' rule,
\[
\begin{aligned}
\Prob(K=k,J=j\mid F)
&=\frac{\Prob(F\mid K=k,J=j)
\Prob(K=k,J=j)}{\Prob(F)}\\
&=\frac{\frac12\Prob(K=k,J=j)}{\frac12}\\
&=\Prob(K=k,J=j).
\end{aligned}
\]

Likewise,
\[
\begin{aligned}
\Prob(K=k,J=j\mid M)
&=\frac{\Prob(M\mid K=k,J=j)
\Prob(K=k,J=j)}{\Prob(M)}\\
&=\frac{\frac12\Prob(K=k,J=j)}{\frac12}\\
&=\Prob(K=k,J=j).
\end{aligned}
\]

Therefore,
\[
\Prob(K=k,J=j\mid F)
=
\Prob(K=k,J=j\mid M)
=
\Prob(K=k,J=j).
\]

Neither system changes the joint distribution of $(K,J)$.

In particular,
\[
\Prob(K=1\mid F)
=
\Prob(K=1\mid M)
=
\frac12.
\]

\medskip
\textbf{Comparison with the original experiment:}

Originally, both rounds shared the same host $H$.

Although $C$ and $J$ were independent conditional on
$(H,K)$, they were generally dependent after averaging over $H$.

For example,
\[
\begin{aligned}
\Prob(C=2,J=2\mid K=1)
&=\frac12
\left[
\left(\frac14\right)^2+
\left(\frac34\right)^2
\right]\\
&=\frac5{16}.
\end{aligned}
\]

But
\[
\begin{aligned}
&\Prob(C=2\mid K=1)\Prob(J=2\mid K=1)\\
&=\frac12\cdot\frac12\\
&=\frac14.
\end{aligned}
\]

\[
\frac5{16}\neq\frac14.
\]

Moreover, in the original experiment,
\[
\begin{aligned}
\Prob(M\mid K=1)
&=\frac12
\left[
\left(\frac14\right)^2+
\left(\frac34\right)^2
\right]\\
&\quad+
\frac12
\left[
\left(\frac34\right)^2+
\left(\frac14\right)^2
\right]\\
&=\frac58,
\end{aligned}
\]
whereas
\[
\Prob(M\mid K=2)
=
\Prob(M\mid K=3)
=
\frac12.
\]

Thus
\[
\Prob(M\mid K=1)
>
\Prob(M\mid K=2)
=
\Prob(M\mid K=3).
\]

The shared host created dependence between the rounds,
so the matching condition $C=J$ favored $K=1$.

Independent host selection removes this dependence and makes
both retention probabilities constant for all $(k,j)$.




\clearpage
\section*{Problem 5}

Start with one amoeba at time $0$. At each minute's end, every living
amoeba dies or splits into two. The death probability is $1/4$ in odd
minutes and $p\in(0,1)$ in even minutes. All choices are independent
across individuals and minutes; an empty population stays empty.
Do not use general branching-process theorems.

\begin{enumerate}
\item For integers $n\geq0$, let $x_n,y_n$ be the probabilities of
extinction within $n$ updates, starting with one amoeba under the
alternating rules, with the odd/even rule first, respectively.
Let $x,y$ be the corresponding eventual probabilities.
Find the initial values, first-step recursions and limiting equations.
Explain why $x=y$ cannot be assumed.
You may use
\mbox{$\Prob(\bigcup_{n\geq0}E_n)
=\lim_{n\to\infty}\Prob(E_n)$}
whenever $E_n\subseteq E_{n+1}$ for all $n\geq0$.

\item Set $p=2/3$. Prove $x=y=1$. An averaged model replaces both
death probabilities by their average $\bar p=11/24$, keeping all
other rules unchanged; its eventual extinction probability is
given as $11/13$.
Calculate extinction probabilities by minute $2$ in both models.
Does averaging overestimate or underestimate extinction at minute $2$
and eventually?

\item Now $p$ is unknown, and the probability of exactly two amoebae
at the end of minute $2$ is $1/3$.
Given exactly four at the end of minute $3$, find every possible
probability that all four descend from the same one of the two
minute-$1$ amoebae, with the corresponding $p$.
Is this conditional probability uniquely determined?
\end{enumerate}

\noindent\textbf{Solution:}

\begin{enumerate}

%-------------------- (a) --------------------

\item[(a)]

Initially, one amoeba is alive:
\[
x_0=y_0=0.
\]

After the first update, extinction occurs either immediately
or when both descendants subsequently become extinct.

\[
x_{n+1}
=\frac14+\frac34 y_n^2,
\]
\[
y_{n+1}
=p+(1-p)x_n^2,
\qquad n\geq0.
\]

In particular,
\[
x_1=\frac14,\qquad y_1=p.
\]

Let $E_n^{(x)}$ and $E_n^{(y)}$ denote extinction within
$n$ updates under the two starting rules.

\[
E_n^{(x)}\subseteq E_{n+1}^{(x)},
\qquad
E_n^{(y)}\subseteq E_{n+1}^{(y)}.
\]

\[
0\leq x_n\leq x_{n+1}\leq1,
\qquad
0\leq y_n\leq y_{n+1}\leq1.
\]

Thus,
\[
x=\lim_{n\to\infty}x_n,
\qquad
y=\lim_{n\to\infty}y_n.
\]

Taking limits in the recursions,
\[
\begin{aligned}
x&=\frac14+\frac34 y^2,\\
y&=p+(1-p)x^2.
\end{aligned}
\]

The two starting rules have different death probabilities:
\[
x_1=\frac14,\qquad y_1=p.
\]

Moreover, if $x=y=q$, then
\[
\frac14+\frac34q^2
=p+(1-p)q^2,
\]
\[
\left(\frac14-p\right)(1-q^2)=0.
\]

\[
\implies
p=\frac14
\quad\text{or}\quad
q=1.
\]

To see that equality need not hold, take $p=1/3$.
Since
\[
\frac14+\frac34\left(\frac12\right)^2
=\frac7{16}<\frac12,
\]
\[
\frac13+\frac23\left(\frac12\right)^2
=\frac12,
\]
induction gives
\[
x_n,y_n\leq\frac12
\qquad(n\geq0).
\]

Hence
\[
x,y\leq\frac12<1,
\qquad p=\frac13\ne\frac14,
\]
\[
\implies x\ne y.
\]

Therefore, $x=y$ cannot be assumed in general.

%-------------------- (b) --------------------

\item[(b)]

For $p=2/3$, the limiting equations become
\[
\begin{cases}
x=\dfrac14+\dfrac34y^2,\\[6pt]
y=\dfrac23+\dfrac13x^2.
\end{cases}
\]

Rearranging,
\[
1-x=\frac34(1-y^2)
=\frac34(1-y)(1+y),
\]
\[
1-y=\frac13(1-x^2)
=\frac13(1-x)(1+x).
\]

If $x=1$, then $y=1$; similarly,
if $y=1$, then $x=1$.

Suppose $x<1$ and $y<1$.

Multiplying the two equations,
\[
(1-x)(1-y)
=\frac14(1-x)(1-y)(1+x)(1+y).
\]

Since $(1-x)(1-y)>0$,
\[
1=\frac14(1+x)(1+y),
\]
\[
(1+x)(1+y)=4.
\]

However,
\[
0\leq x,y<1
\implies
(1+x)(1+y)<4,
\]
a contradiction.

Thus,
\[
x=y=1.
\]

\medskip
\noindent\textbf{Extinction by minute 2:}

For the original alternating model,
\[
\begin{aligned}
P_{\mathrm{alt}}(T\leq2)
&=x_2\\
&=\frac14+\frac34y_1^2\\
&=\frac14+\frac34\left(\frac23\right)^2\\
&=\frac14+\frac13\\
&=\frac7{12}.
\end{aligned}
\]

For the averaged model,
\[
\bar p
=\frac12\left(\frac14+\frac23\right)
=\frac{11}{24},
\]
\[
1-\bar p=\frac{13}{24}.
\]

\[
\begin{aligned}
P_{\mathrm{avg}}(T\leq2)
&=\bar p+(1-\bar p)\bar p^2\\
&=\frac{11}{24}
+\frac{13}{24}\left(\frac{11}{24}\right)^2\\
&=\frac{11}{24}+\frac{1573}{13824}\\
&=\frac{7909}{13824}.
\end{aligned}
\]

Comparison:
\[
\frac7{12}
=\frac{8064}{13824}
>\frac{7909}{13824}.
\]

\[
P_{\mathrm{alt}}(T\leq2)
-P_{\mathrm{avg}}(T\leq2)
=\frac{155}{13824}>0.
\]

Thus, averaging underestimates extinction by minute $2$.

\medskip
\noindent\textbf{Eventual extinction:}

\[
P_{\mathrm{alt}}(T<\infty)=1,
\]
\[
P_{\mathrm{avg}}(T<\infty)=\frac{11}{13}.
\]

\[
1-\frac{11}{13}=\frac2{13}>0.
\]

Thus, averaging also underestimates eventual extinction.

%-------------------- (c) --------------------

\item[(c)]

Let $N_k$ denote the number of amoebae at the end
of minute $k$.

At minute $1$,
\[
P(N_1=2)=\frac34.
\]

Exactly two amoebae at minute $2$ requires exactly
one of the two minute-$1$ amoebae to split.

\[
P(N_2=2)
=\frac34\binom21(1-p)p.
\]

Given $P(N_2=2)=1/3$,
\[
\frac32p(1-p)=\frac13,
\]
\[
p(1-p)=\frac29,
\]
\[
p^2-p+\frac29=0,
\]
\[
(3p-1)(3p-2)=0.
\]

Thus,
\[
p=\frac13
\quad\text{or}\quad
p=\frac23.
\]

Let
\[
A=\{\text{all four minute-3 amoebae descend
from the same minute-1 amoeba}\},
\]
\[
B=\{N_3=4\}.
\]

Since $B\subseteq\{N_1=2\}$,
\[
P(A\mid B)
=\frac{P(A\cap B\mid N_1=2)}
{P(B\mid N_1=2)}.
\]

There are two possible cases.

\medskip
\noindent\textbf{Case 1: $N_2=2$.}

Exactly one minute-$1$ amoeba splits.
Both minute-$2$ amoebae must split at minute $3$.

\[
P(N_2=2\mid N_1=2)
=2p(1-p).
\]

\[
P(B\mid N_2=2)
=\left(\frac34\right)^2.
\]

All four descendants necessarily have the same
minute-$1$ ancestor.

\medskip
\noindent\textbf{Case 2: $N_2=4$.}

Both minute-$1$ amoebae split:
\[
P(N_2=4\mid N_1=2)=(1-p)^2.
\]

For exactly four amoebae at minute $3$,
exactly two of the four minute-$2$ amoebae must split.

\[
P(B\mid N_2=4)
=\binom42
\left(\frac34\right)^2
\left(\frac14\right)^2.
\]

For all four to have the same minute-$1$ ancestor,
the two splitting amoebae must be the two descendants
of the same minute-$1$ amoeba.

There are $2$ favorable choices among
$\binom42=6$ possible pairs.

\[
P(A\cap B\mid N_2=4)
=2\left(\frac34\right)^2
\left(\frac14\right)^2.
\]

Combining both cases,
\[
\begin{aligned}
P(B\mid N_1=2)
={}&2p(1-p)\left(\frac34\right)^2\\
&+(1-p)^2\binom42
\left(\frac34\right)^2
\left(\frac14\right)^2.
\end{aligned}
\]

\[
\begin{aligned}
P(A\cap B\mid N_1=2)
={}&2p(1-p)\left(\frac34\right)^2\\
&+2(1-p)^2
\left(\frac34\right)^2
\left(\frac14\right)^2.
\end{aligned}
\]

Therefore,
\[
\begin{aligned}
P(A\mid B)
&=\frac{
2p(1-p)+\dfrac18(1-p)^2
}{
2p(1-p)+\dfrac38(1-p)^2
}\\[6pt]
&=\frac{
2p+\dfrac18(1-p)
}{
2p+\dfrac38(1-p)
}\\[6pt]
&=\frac{16p+1-p}{16p+3(1-p)}\\[6pt]
&=\frac{1+15p}{3+13p}.
\end{aligned}
\]

For $p=1/3$,
\[
\begin{aligned}
P(A\mid B)
&=\frac{1+15(1/3)}{3+13(1/3)}\\
&=\frac6{22/3}\\
&=\frac9{11}.
\end{aligned}
\]

For $p=2/3$,
\[
\begin{aligned}
P(A\mid B)
&=\frac{1+15(2/3)}{3+13(2/3)}\\
&=\frac{11}{35/3}\\
&=\frac{33}{35}.
\end{aligned}
\]

The two possible pairs are
\[
\begin{aligned}
p=\frac13
&\quad\Longrightarrow\quad P(A\mid B)=\frac9{11},\\[4pt]
p=\frac23
&\quad\Longrightarrow\quad P(A\mid B)=\frac{33}{35}.
\end{aligned}
\]

Since
\[
\frac9{11}\ne\frac{33}{35},
\]
the conditional probability is not uniquely determined.

\end{enumerate}

\end{document}